\section{Functional Access Control}

Functional access control is the core \SCS{} authorization model. Callers normally do not pass capability tokens. Methods declare required authority. The enforcement module derives the required authority from:

\begin{lstlisting}
contract_id
method_name
method_arguments
msg_sender
tx_sender
current state
current registry
\end{lstlisting}

\subsection{Default call shape}

Preferred:

\begin{lstlisting}
(call-contract! tx TokenA transferFrom Alice Bob USDC 10)
\end{lstlisting}

Avoid by default:

\begin{lstlisting}
(call-contract! tx TokenA transferFrom (CapRef Cap-123) Alice Bob USDC 10)
\end{lstlisting}

Capability references \MAY{} appear in specialized methods, but they are not the ordinary programming model.

\subsection{Method policy structure}

Conceptual policy form:

\begin{lstlisting}
(method-policy TokenA transferFrom
  (principal msg_sender)

  (requires
    AuthenticatedCaller
    ContractNotPaused
    (RegistryGrant
      (registry TokenA.capabilities)
      (holder msg_sender)
      (resource (Allowance TokenA $from $asset))
      (right transfer-from)
      (amount $amount)
      (consume $amount)))

  (reads
    (Balance TokenA $from $asset)
    (Balance TokenA $to $asset)
    (Allowance TokenA $from msg_sender $asset))

  (writes
    (Balance TokenA $from $asset)
    (Balance TokenA $to $asset)
    (Allowance TokenA $from msg_sender $asset))

  (emits Transfer)

  (invariants
    NoNegativeBalances
    TotalSupplyConserved)

  (reentrancy nonreentrant))
\end{lstlisting}

\subsection{Functional access-control requirements}

\req{SCS-FAC-001 --- Policy per public method}{Every public state-changing method \MUST{} have a registered method policy.}

\req{SCS-FAC-002 --- Deny by default}{A method with no registered policy \MUST{} be unavailable to external callers.}

\req{SCS-FAC-003 --- Runtime-derived principal}{The authorization principal \MUST{} be derived from transaction context, not from a user-provided method argument.}

\req{SCS-FAC-004 --- Default principal}{The default authorization principal \SHOULD{} be \code{msg\_sender}.}

\req{SCS-FAC-005 --- Live registry lookup}{The enforcement module \MUST{} check required durable authority through live capability-registry lookup.}

\req{SCS-FAC-006 --- No routine capability arguments}{Normal methods \SHOULDNOT{} require callers to pass capability tokens or capability references.}

\req{SCS-FAC-007 --- Argument-dependent policies}{Method policies \MAY{} bind requirements to method arguments. For example, \code{transferFrom(from, to, asset, amount)} requires an allowance grant from \code{from} to \code{msg\_sender} for \code{asset} and at least \code{amount}.}

\req{SCS-FAC-008 --- Guarded policy registration}{Only deployment, governance, or authorized admin mechanisms \MAY{} register or update method policies.}

\req{SCS-FAC-009 --- Dispatcher enforcement}{The dispatcher \MUST{} enforce method policy before private implementation execution.}

\req{SCS-FAC-010 --- Private implementation protection}{The private implementation behind a public method \MUSTNOT{} be externally callable.}

\req{SCS-FAC-011 --- Atomic authorization effects}{Any registry effects caused by authorization, such as allowance consumption or nonce consumption, \MUST{} commit or revert atomically with the contract call.}

\req{SCS-FAC-012 --- Method-scoped write set}{A method policy \SHOULD{} declare the state classes or keys the method may write.}

\req{SCS-FAC-013 --- Storage-level write-scope enforcement}{The storage kernel \MUST{} reject writes outside the active method's authorized write scope.}

\req{SCS-FAC-014 --- View policies}{Read-only methods \SHOULD{} have policies too, even if the policy is simply \code{PublicRead}.}

\req{SCS-FAC-015 --- Auditability}{Authorization success, authorization failure, and registry consumption \SHOULD{} be auditable.}

\subsection{Two-layer authorization}

\SCS{} separates method access from state effect control:

\begin{lstlisting}
method access control:
    who may call this method?

effect control:
    what may this method mutate?
\end{lstlisting}

Both layers are required. A caller may be authorized to call \code{transferFrom}, but the implementation must still be prevented from mutating unrelated state such as minter roles, upgrade admin records, or total supply unless the policy explicitly permits those writes.
